% GENERATED FILE. Edit canonical lessons, then run npm run generate:books.
% canonical-source-hash: fnv1a64:90845aa5f84cb894
% canonical-lessons: GU-C11-mitra, GU-C11-kutumb, GU-C11-bhai, GU-C11-bahen

\chapter{Friend and Family}
\label{ch:gu-friend-and-family}
\hlchaptermodality{11}

\begin{chapteropening}
\textbf{By the end of this chapter:} \emph{I can name a friend, a family, a brother, and a sister in Gujarati, explain which of these words was inherited whole from Sanskrit rather than worn down by sound change, and say which sibling word is genuinely a cousin of its English counterpart and which is not.}

The chapter closes on \gu{બહેન}, deliberately not a cousin of English sister: the reader produces \gu{મિત્ર}, \gu{કુટુંબ}, \gu{ભાઈ} and \gu{બહેન}, names \gu{ભાઈ} as English brother's unbroken cousin against \gu{બહેન}'s unrelated (disputed bhaga-based) root, and names \gu{કુટુંબ} as taken up whole from Sanskrit rather than worn down like the other three — while recalling that \gu{મિત્ર} and \gu{મહેરબાની}, taught a chapter apart, share the same ancient Indo-Iranian root.
\end{chapteropening}

\section[mitra]{\gu{મિત્ર} — a friend, and the word inside your own \textquotedblleft{}please\textquotedblright{}}

\label{lesson:GU-C11-mitra}



\begin{quote}
Last chapter asked politely for \textbf{roṭlī} with \emph{meherbānī karīne}. That phrase is about to turn out to share its root with the very first word of this new chapter.
\end{quote}

\subsection*{You'll want to know first — one word}
\begin{quote}
\textbf{\gu{મિત્ર}} — \emph{mitra} — \textbf{friend}
\end{quote}

\subsection*{The letters in this word}
\textbf{\gu{મ}} \emph{ma}, \textbf{\gu{િ}} the short-\emph{i} sign written \textbf{before} the consonant it belongs to — the same leftward quirk you saw open \emph{vichārvũ} — then a conjunct: \textbf{\gu{ત્ર}} — \textbf{\gu{ત}} \emph{ta} with the virama \textbf{\gu{્}} stripping its vowel, joined straight into \textbf{\gu{ર}} \emph{ra}. \emph{Mi-tra}, two syllables, the second one a cluster.

\begin{cousinweb}
\textbf{mitra} is Sanskrit for \textquotedblleft{}friend,\textquotedblright{} from Proto-Indo-Iranian \emph{\textbf{mitras}, \textquotedblleft{}(that which) causes binding\textquotedblright{} — a covenant word before it was a friendship word. Its exact twin is Avestan \textbf{Mithra}, the god of contracts and given word, whose name lived on in Middle Persian as \textbf{Mihr}.}

That is not a dead end. \textbf{\gu{મહેરબાની}}, the \textquotedblleft{}kindness\textquotedblright{} you built \emph{meherbānī karīne} from last chapter, is Persian \emph{mihrbānī} — built on that very same \textbf{Mihr}. So \textbf{\gu{મિત્ર}}, inherited straight down from Sanskrit, and \textbf{\gu{મહેરબાની}}, borrowed centuries later from Persian, are cousins by the same Indo-Iranian root — one word for \textquotedblleft{}friend,\textquotedblright{} the other for \textquotedblleft{}favor,\textquotedblright{} both once naming the same idea of a bond that holds.
\end{cousinweb}

\subsection*{Your turn}
\begin{itemize}
\raggedright
  \item \emph{Say it:} \textbf{mitra} — friend
  \item \emph{Trace it:} Proto-Indo-Iranian \emph{mitras} $\to$ Avestan \textbf{Mithra} $\to$ Persian \textbf{Mihr} $\to$ \textbf{\gu{મહેરબાની}}
  \item \emph{Connect:} \textbf{\gu{મિત્ર}} and \textbf{\gu{મહેરબાની}} — same ancient root, two different doors into Gujarati
  \item \emph{Run through:} \textbf{roṭlī, meherbānī karīne} once more, now naming the \textquotedblleft{}friend\textquotedblright{} root inside it
\end{itemize}

\begin{tcolorbox}[breakable,colback=teal!4,colframe=teal!35!black,title={Before you move on}]
Give the word for \textquotedblleft{}friend.\textquotedblright{} (\emph{Mitra}.) Which other Gujarati word you already know shares its root, and through which language did each arrive? (\textbf{\gu{મહેરબાની}} — \emph{mitra} inherited from Sanskrit, \emph{meherbānī} borrowed from Persian \emph{mihr}.) What did the shared root originally mean? (\textquotedblleft{}\textbf{To bind},\textquotedblright{} a covenant.)

Sources: \href{https://en.wiktionary.org/wiki/\%E0\%AA\%AE\%E0\%AA\%BF\%E0\%AA\%A4\%E0\%AB\%8D\%E0\%AA\%B0}{Wiktionary: \gu{મિત્ર}}; \href{https://en.wiktionary.org/wiki/Mithra}{Wiktionary: Mithra}.
\end{tcolorbox}
\section[kuṭumb]{\gu{કુટુંબ} — a word that skipped the sound changes}

\label{lesson:GU-C11-kutumb}



\begin{quote}
\textbf{\gu{મિત્ર}} (\emph{mitra}) reached Gujarati the long way, worn down through Prakrit like nearly every word you own. The word for the group a friend gets folded into took a different road entirely.
\end{quote}

\subsection*{You'll want to know first — one word}
\begin{quote}
\textbf{\gu{કુટુંબ}} — \emph{kuṭumb} — \textbf{family}
\end{quote}

\subsection*{The letters in this word}
\textbf{\gu{ક}} \emph{ka}, \textbf{\gu{ુ}} the short-\emph{u} sign below it, then \textbf{\gu{ટ}} — the same retroflex \emph{ṭ} you first curled back in \emph{roṭlī}, here doing the same job again. \textbf{\gu{ું}} closes it: \emph{u} with the nasal dot on top, then \textbf{\gu{બ}} \emph{ba}. \emph{Ku-ṭu-m-b}.

\begin{grammarlens}[title={Grammar Lens: taken whole, not worn down}]
You have seen Sanskrit reach back into Gujarati once before: \textbf{\gu{ત્રણ}}'s \emph{r}, restored after Prakrit had already dropped it — a \textbf{repair}, made to a word that had kept changing the whole time. \textbf{\gu{કુટુંબ}} is a different kind of contact. Gujarati did not inherit this word through the centuries of sound change that gave you \emph{pāṇī} or \emph{be}; it took the Sanskrit word up \textbf{directly}, close to whole, the way a modern language reaches for a ready-made technical term instead of growing one on its own. Linguists call an inherited word like \emph{pāṇī} \textbf{tadbhava}, \textquotedblleft{}born from that,\textquotedblright{} and a direct borrowing like this one \textbf{tatsama}, \textquotedblleft{}same as that\textquotedblright{} — \emph{kuṭumb} is tatsama.
\end{grammarlens}

\begin{cousinweb}
Here the trail runs further than usual, in an unexpected direction. Sanskrit \textbf{kuṭumba} is itself thought to be a \textbf{loanword}, taken from a Dravidian source and compared to Tamil \emph{kuṭimai}, \textquotedblleft{}family, lineage,\textquotedblright{} both built on \emph{kuṭi}, \textquotedblleft{}hut.\textquotedblright{} A word that looks purely Sanskrit, and that Gujarati then borrowed as purely Sanskrit, may have entered Sanskrit from a completely different family of languages first. The chain of ownership runs Dravidian $\to$ Sanskrit $\to$ Gujarati, with no step skipped and no step certain beyond reasonable comparison.
\end{cousinweb}

\subsection*{Your turn}
\begin{itemize}
\raggedright
  \item \emph{Say it:} \textbf{kuṭumb} — family
  \item \emph{Contrast:} \textbf{tadbhava} \emph{pāṇī} (worn down over time) against \textbf{tatsama} \emph{kuṭumb} (taken whole)
  \item \emph{Match:} this \textquotedblleft{}taken whole\textquotedblright{} idea against \emph{traṇ}'s restored \emph{r} — both are Sanskrit reaching forward, in different ways
  \item \emph{Trace it:} proposed Dravidian \emph{kuṭi} \textquotedblleft{}hut\textquotedblright{} $\to$ Sanskrit \emph{kuṭumba} $\to$ \textbf{\gu{કુટુંબ}}
\end{itemize}

\begin{tcolorbox}[breakable,colback=teal!4,colframe=teal!35!black,title={Before you move on}]
Give the word for \textquotedblleft{}family.\textquotedblright{} (\emph{Kuṭumb}.) Is it \emph{tadbhava} or \emph{tatsama}? (\textbf{Tatsama} — borrowed whole from Sanskrit, not worn down through Prakrit.) Where does Sanskrit's own \emph{kuṭumba} likely come from? (\textbf{Dravidian}, compared to Tamil \emph{kuṭimai}.)

Sources: \href{https://en.wiktionary.org/wiki/\%E0\%AA\%95\%E0\%AB\%81\%E0\%AA\%9F\%E0\%AB\%81\%E0\%AA\%82\%E0\%AA\%AC}{Wiktionary: \gu{કુટુંબ}}; \href{https://en.wiktionary.org/wiki/\%E0\%A4\%95\%E0\%A5\%81\%E0\%A4\%9F\%E0\%A5\%81\%E0\%A4\%AE\%E0\%A5\%8D\%E0\%A4\%AC}{Wiktionary: Sanskrit kutumba}.
\end{tcolorbox}
\section[bhāī]{\gu{ભાઈ} — the same word as \textquotedblleft{}brother,\textquotedblright{} unbroken}

\label{lesson:GU-C11-bhai}



\begin{quote}
\textbf{\gu{કુટુંબ}} (\emph{kuṭumb}) was taken up whole, close to the finished Sanskrit shape. The two words for its members run the opposite way — worn all the way down, the ordinary path nearly every Gujarati word has followed.
\end{quote}

\subsection*{You'll want to know first — one word}
\begin{quote}
\textbf{\gu{ભાઈ}} — \emph{bhāī} — \textbf{brother}
\end{quote}

\subsection*{The letters in this word}
\textbf{\gu{ભ}} \emph{bha}, the aspirated \emph{b}, with the long-\emph{ā} sign \textbf{\gu{ા}}, then \textbf{\gu{ઈ}}, the independent long-\emph{ī} vowel standing on its own rather than attached to a consonant. \emph{Bhā-ī}, two syllables, no consonant closing the second.

\begin{cousinweb}
\textbf{bhāī} is inherited, through Prakrit \emph{bhāi}, from Sanskrit \textbf{bhrātṛ}. Its root is Indo-European *\emph{b\textsuperscript{h}réh\textsubscript{2}tēr}, and this is one of the cleanest cognate sets Indo-European linguistics has: Sanskrit \emph{bhrātṛ}, English \textbf{brother}, Latin \emph{frāter}, Greek \emph{phrātēr}, Russian \emph{brat} — every one of them the same word, carried down by ordinary sound change in every branch rather than borrowed by any of them. Where \emph{kuṭumb} skipped the centuries of change that shape \emph{tadbhava} words, \emph{bhāī} is the textbook case of what those centuries do: the consonant cluster \emph{bhr-} simplified once the \emph{r} dropped, exactly the kind of wearing-down that gave you \emph{pāṇī} from \emph{pānīya}.
\end{cousinweb}

\subsection*{Your turn}
\begin{itemize}
\raggedright
  \item \emph{Say it:} \textbf{bhāī} — brother
  \item \emph{Trace it:} PIE \emph{b\textsuperscript{h}réh\textsubscript{2}tēr} $\to$ Sanskrit \emph{bhrātṛ} $\to$ \textbf{\gu{ભાઈ}}, and $\to$ English \textbf{brother}
  \item \emph{Contrast:} \emph{bhāī} worn down over centuries against \emph{kuṭumb} taken up whole
  \item \emph{Name:} two other languages whose word for \textquotedblleft{}brother\textquotedblright{} is the same cognate (Latin \emph{frāter}, Greek \emph{phrātēr}, or Russian \emph{brat})
\end{itemize}

\begin{tcolorbox}[breakable,colback=teal!4,colframe=teal!35!black,title={Before you move on}]
Give the word for \textquotedblleft{}brother.\textquotedblright{} (\emph{Bhāī}.) Which English word is the very same word, by unbroken descent, not borrowing? (\textbf{Brother}.) Is \emph{bhāī} \emph{tadbhava} or \emph{tatsama} — worn down, or taken whole like \emph{kuṭumb}? (\textbf{Tadbhava} — worn down.)

Sources: \href{https://en.wiktionary.org/wiki/\%E0\%AA\%AD\%E0\%AA\%BE\%E0\%AA\%88}{Wiktionary: \gu{ભાઈ}}; \href{https://en.wiktionary.org/wiki/Reconstruction:Proto-Indo-European/b\%CA\%B0r\%C3\%A9h\%E2\%82\%82t\%C4\%93r}{Wiktionary: b\textsuperscript{h}réh\textsubscript{2}tēr}.
\end{tcolorbox}
\section[bahen]{\gu{બહેન} — the sister who is not \textquotedblleft{}sister\textquotedblright{}'s cousin}

\label{lesson:GU-C11-bahen}



\begin{quote}
\textbf{\gu{ભાઈ}} is English \textbf{brother}, unbroken, down to the root. Do not assume the sibling pair matches. This chapter's last word breaks the pattern.
\end{quote}

\subsection*{You'll want to know first — one word}
\begin{quote}
\textbf{\gu{બહેન}} — \emph{bahen} — \textbf{sister}
\end{quote}

\subsection*{The letters in this word}
\textbf{\gu{બ}} \emph{ba}, \textbf{\gu{હ}} \emph{ha}, \textbf{\gu{ે}} the \emph{e} sign above the consonant, then \textbf{\gu{ન}} \emph{na}. \emph{Ba-hen}, every letter already yours from earlier chapters.

\begin{cousinweb}
Here is the surprise: \textbf{\gu{બહેન}} is \textbf{not} \emph{bhāī}'s counterpart the way English \emph{sister} is \emph{brother}'s. English \emph{sister}, Latin \emph{soror}, and Sanskrit \textbf{svasṛ} all descend from one shared PIE root, *\emph{swésōr}. Gujarati's word for sister does not use that root at all.

\textbf{bahen} comes from Sanskrit \textbf{bhaginī}, through Prakrit \emph{bhaginī}. Scholars connect \emph{bhaginī} to \textbf{bhaga}, \textquotedblleft{}good fortune, a share\textquotedblright{} — the same \emph{bhaga} behind \emph{bhāgya}, \textquotedblleft{}fate\textquotedblright{} — plus a feminine suffix, so \textquotedblleft{}sister\textquotedblright{} may originally have meant something closer to \textquotedblleft{}\textbf{the one who shares}.\textquotedblright{} That connection is debated rather than settled, the way \emph{roṭlī}'s root was: some scholars accept the \emph{bhaga} link, others have argued against it. What is certain is what it is \textbf{not}: unlike \emph{bhāī}, \emph{bahen} is no relation to English \emph{sister} at all — and unlike the disputed-but-real \emph{mitra}/\emph{meherbānī} pair, this dial lands on \textquotedblleft{}unrelated,\textquotedblright{} full stop.
\end{cousinweb}

\subsection*{Your turn}
\begin{itemize}
\raggedright
  \item \emph{Say it:} \textbf{bahen} — sister, and \textbf{bhāī} — brother
  \item \emph{Contrast:} \emph{bhāī} — same root as English \emph{brother}; \emph{bahen} — a different Sanskrit word entirely, not \emph{sister}'s cousin
  \item \emph{Trace it:} proposed \emph{bhaga} \textquotedblleft{}share, fortune\textquotedblright{} $\to$ \emph{bhaginī} $\to$ \textbf{\gu{બહેન}}
  \item \emph{Run through:} all four family words — \textbf{mitra}, \textbf{kuṭumb}, \textbf{bhāī}, \textbf{bahen} — inherited/borrowed, direct/indirect
  \item \emph{Recall:} \textbf{\gu{મિત્ર}} and \textbf{\gu{મહેરબાની}}'s shared root from the last chapter, and \textbf{\gu{કુટુંબ}}'s proposed Dravidian root
\end{itemize}

\begin{tcolorbox}[breakable,colback=teal!4,colframe=teal!35!black,title={Before you move on}]
Give the word for \textquotedblleft{}sister.\textquotedblright{} (\emph{Bahen}.) Is it the same root as English \emph{sister}? (\textbf{No} — that connection belongs to \emph{bhāī}/\textbf{brother} instead.) What does \emph{bahen} probably come from? (Sanskrit \emph{bhaginī}, disputedly tied to \emph{bhaga}, \textquotedblleft{}share.\textquotedblright{}) Name this chapter's four people and things: friend, family, brother, sister. (\emph{Mitra}, \emph{kuṭumb}, \emph{bhāī}, \emph{bahen}.)

Sources: \href{https://en.wiktionary.org/wiki/\%E0\%AA\%AC\%E0\%AA\%B9\%E0\%AB\%87\%E0\%AA\%A8}{Wiktionary: \gu{બહેન}}; \href{https://en.wiktionary.org/wiki/\%E0\%A4\%AD\%E0\%A4\%97\%E0\%A4\%BF\%E0\%A4\%A8\%E0\%A5\%80}{Wiktionary: Sanskrit bhagini}.
\end{tcolorbox}
